% Generated draft: requires mathematical and source review.
\documentclass[11pt]{article}
\usepackage[margin=25mm]{geometry}
\usepackage{fontspec}
\setmainfont{Latin Modern Roman}
\usepackage{amsmath,amssymb,mathtools,unicode-math}
\setmathfont{Latin Modern Math}
\usepackage{enumitem,needspace,xurl,hyperref}
\setlength{\parindent}{0pt}
\setlength{\parskip}{5pt}
\setlength{\emergencystretch}{3em}
\newcommand{\N}{\mathbb{N}}
\newcommand{\Z}{\mathbb{Z}}
\newcommand{\Q}{\mathbb{Q}}
\newcommand{\R}{\mathbb{R}}
\newcommand{\C}{\mathbb{C}}
\newcommand{\F}{\mathbb{F}}
\newcommand{\A}{\mathbb{A}}
\newcommand{\PP}{\mathbb{P}}
\newcommand{\E}{\mathbb{E}}
\newcommand{\eps}{\varepsilon}
\newcommand{\dd}{\,\mathrm{d}}
\begin{document}
\section*{UnsolvedMath ID: 30003307}
\textbf{Source problem code:} catalogue entry
\subsection{Definitions and mathematical statement}
For a Noetherian local ring $(A,\mathfrak m)$ of dimension $d$, its Hilbert--Samuel multiplicity is
\[
e(A)=\lim_{n\to\infty}\frac{d!}{n^d}\operatorname{length}_A(A/\mathfrak m^{n+1}).
\]
Write $e_X(x)=e(\mathcal O_{X,x})$ when this definition applies. This is the usual interpretation of multiplicity adopted for explaining the question; the supplied record does not explicitly specify a multiplicity invariant. A closed center $C\subseteq X$ is regular if its local rings are regular. Under the usual multiplicity interpretation, equimultiplicity along $C$ means that $e_X(c)$ is constant for $c\in C$. Whether the intended centers must additionally lie in the locus of maximum multiplicity is not specified. If $\mathcal I_C$ is the ideal sheaf of $C$, the blowup along $C$ is
\[
\operatorname{Bl}_C(X)=\operatorname{Proj}_X\!\left(\bigoplus_{n\geq0}\mathcal I_C^n\right).
\]
\textbf{Source question.} Can the multiplicity of a variety in positive characteristic be reduced by successive blowups along regular equimultiple centers? In particular, can such reduction be proved in the situation
\[
X=\operatorname{Spec}(B)\longrightarrow W=\operatorname{Spec}(S),
\qquad S\subset B,\qquad S\text{ regular},
\]
where the inclusion $S\subset B$ is described in the source as purely inseparable?

The record does not specify the base field, the precise ring-theoretic meaning of purely inseparable, or finiteness hypotheses on $S\subset B$. It also does not specify whether reduction is local at a chosen point, concerns maximum multiplicity globally, or requires a strict decrease at every point above a chosen point after finitely many blowups. Thus the question above preserves the supplied target, but these missing qualifications prevent a unique fully quantified conjecture. The record also leaves unspecified whether the blowups are intrinsic to $X$ or are performed in an ambient space with corresponding transforms.
\subsection{Short English statement}
The question asks whether blowing up regular centers along which multiplicity is constant can lower multiplicity in positive characteristic, especially for a purely inseparable affine morphism to a regular affine scheme. The exact hypotheses and the required meaning of reduction remain unspecified. This is a statement-only draft requiring human mathematical review.
\subsection{Sources}
\begin{itemize}
\item Local registry: \texttt{data/unsolvedmath/problems.json}, supplied record with permanent UnsolvedMath ID \texttt{30003307}, source problem code \texttt{catalogue entry}, title ``Multiplicity Reduction by Regular Equimultiple Blowups.'' No individual proposer is supplied; the proposed year is recorded as 2016.
\item Supplied report citation: \emph{Asymptotic Phenomena in Local Algebra and Singularity Theory} (2017), Oberwolfach Reports, pp.\ 3199--3259, DOI \texttt{10.4171/owr/2016/57}. \url{https://doi.org/10.4171/owr/2016/57}. The background attributes the question to workshop open-problem discussions collected in Oberwolfach Reports; no precise problem page or named speaker is supplied.
\item Supplied literature source: ``Hauser, singularities in positive characteristic.'' \url{https://arxiv.org/abs/1711.03453}.
\item Supplied literature source: ``OWR 57 (2017).'' \url{https://doi.org/10.4171/owr/2016/57}.
\item Literature assessment dated exactly \texttt{2026-08-21}: the supplied curation labels the problem open and describes reduction by regular equimultiple blowups as a recognized route to positive-characteristic resolution, with the general reduction problem remaining open. This is supplied dated curation, not an independently verified literature assessment. The background expressly requests independent verification of time-sensitive status and cited claims.
\item The registry labels the statement \texttt{corrected\_verified}; its accompanying qualification says that only notation, typography, and TeX-delimiter cleanup was applied, preserving hypotheses, quantifiers, and subquestions. This label does not establish mathematical correctness or resolve the missing formulation details.
\end{itemize}
\end{document}
